\chapter{MLIR Dialect Developer Quickstart Guide}
\label{app:mlir_quickstart}

\begin{quote}
\textit{``Writing an MLIR dialect is the modern equivalent of writing an instruction set architecture: you define the fundamental types, operations, invariants, and transformation passes that shape computation.''}
\end{quote}

This appendix serves as a practical, hands-on developer manual for setting up, building, and extending the \texttt{dqc} MLIR dialect from source.

% ==============================================================================
\section{Project Directory Structure}
\label{sec:app_project_structure}
% ==============================================================================

An out-of-tree MLIR dialect repository follows standard LLVM project conventions:

\begin{lstlisting}[language=MLIR, caption={Standard Out-of-Tree MLIR Dialect Directory Hierarchy}]
dqc-compiler/
|-- CMakeLists.txt
|-- include/
|   \-- dqc/
|       |-- DQCDialect.h
|       |-- DQCDialect.td          # Root Dialect TableGen
|       |-- DQCTypes.h
|       |-- DQCTypes.td            # Custom Types TableGen
|       |-- DQCOps.h
|       |-- DQCOps.td              # Operation Definitions TableGen
|       \-- Passes/
|           |-- Passes.h
|           \-- Passes.td          # Pass Definitions TableGen
|-- lib/
|   \-- DQC/
|       |-- CMakeLists.txt
|       |-- DQCDialect.cpp         # Dialect Initialization
|       |-- DQCTypes.cpp           # Custom Type Parsers/Printers
|       |-- DQCOps.cpp             # Op Verifiers & Custom Logic
|       \-- Transforms/
|           |-- CanonicalizationPass.cpp
|           |-- HypergraphPartitioningPass.cpp
|           |-- TeleportationSynthesisPass.cpp
|           |-- CoherenceAwareSchedulingPass.cpp
|           \-- LowerToQIRPass.cpp
|-- python/
|   \-- dqc/
|       |-- CMakeLists.txt
|       |-- DQCModule.cpp          # pybind11 C-API Bindings
|       \-- __init__.py
|-- test/
|   |-- CMakeLists.txt
|   |-- lit.cfg.py
|   \-- DQC/
|       |-- pass1_canonicalize.mlir
|       |-- pass2_partition.mlir
|       |-- pass3_teleport.mlir
|       |-- pass4_schedule.mlir
|       \-- pass5_lowering.mlir
\-- tools/
    \-- dqc-opt/
        |-- CMakeLists.txt
        \-- dqc-opt.cpp            # Custom Compiler Driver Binary
\end{lstlisting}

% ==============================================================================
\section{CMake Build System Configuration}
\label{sec:app_cmake_config}
% ==============================================================================

To build against an existing LLVM/MLIR installation:

\begin{lstlisting}[language=C++, caption={Root CMakeLists.txt for DQC Dialect}]
cmake_minimum_required(VERSION 3.20)
project(dqc-compiler LANGUAGES CXX C)

set(CMAKE_CXX_STANDARD 17)
set(CMAKE_CXX_STANDARD_REQUIRED ON)

find_package(MLIR REQUIRED CONFIG)

message(STATUS "Using MLIRConfig.cmake in: ${MLIR_DIR}")
message(STATUS "Using LLVMConfig.cmake in: ${LLVM_DIR}")

set(LLVM_RUNTIME_OUTPUT_INTDIR ${CMAKE_BINARY_DIR}/bin)
set(LLVM_LIBRARY_OUTPUT_INTDIR ${CMAKE_BINARY_DIR}/lib)
set(MLIR_BINARY_DIR ${CMAKE_BINARY_DIR})

list(APPEND CMAKE_MODULE_PATH "${MLIR_CMAKE_DIR}")
list(APPEND CMAKE_MODULE_PATH "${LLVM_CMAKE_DIR}")

include(TableGen)
include(AddLLVM)
include(AddMLIR)
include(HandleLLVMOptions)

include_directories(${LLVM_INCLUDE_DIRS})
include_directories(${MLIR_INCLUDE_DIRS})
include_directories(${PROJECT_SOURCE_DIR}/include)
include_directories(${PROJECT_BINARY_DIR}/include)

add_subdirectory(include/dqc)
add_subdirectory(lib)
add_subdirectory(tools)
add_subdirectory(test)
\end{lstlisting}

\begin{lstlisting}[language=C++, caption={TableGen Include CMakeLists.txt (include/dqc/CMakeLists.txt)}]
set(LLVM_TARGET_DEFINITIONS DQCDialect.td)
mlir_tablegen(DQCDialect.h.inc -gen-dialect-decls)
mlir_tablegen(DQCDialect.cpp.inc -gen-dialect-defs)
add_public_tablegen_target(DQCGenDialect)

set(LLVM_TARGET_DEFINITIONS DQCOps.td)
mlir_tablegen(DQCOps.h.inc -gen-op-decls)
mlir_tablegen(DQCOps.cpp.inc -gen-op-defs)
add_public_tablegen_target(DQCGenOps)

set(LLVM_TARGET_DEFINITIONS DQCTypes.td)
mlir_tablegen(DQCTypes.h.inc -gen-typedef-decls)
mlir_tablegen(DQCTypes.cpp.inc -gen-typedef-defs)
add_public_tablegen_target(DQCGenTypes)
\end{lstlisting}

% ==============================================================================
\section{Driver Implementation: \texttt{dqc-opt.cpp}}
\label{sec:app_dqc_opt}
% ==============================================================================

\begin{lstlisting}[language=C++, caption={Custom Tool Driver (tools/dqc-opt/dqc-opt.cpp)}]
#include "mlir/IR/DialectRegistry.h"
#include "mlir/InitAllDialects.h"
#include "mlir/InitAllPasses.h"
#include "mlir/Tools/mlir-opt/MlirOptMain.h"
#include "dqc/DQCDialect.h"
#include "dqc/Passes/Passes.h"

int main(int argc, char **argv) {
  mlir::DialectRegistry registry;

  // Register Standard MLIR Dialects
  mlir::registerAllDialects(registry);
  mlir::registerAllPasses();

  // Register DQC Custom Dialect & Passes
  registry.insert<mlir::dqc::DQCDialect>();
  mlir::dqc::registerDQCPasses();

  return mlir::asMainReturnCode(
      mlir::MlirOptMain(argc, argv, "Distributed Quantum Compiler Driver\n", registry));
}
\end{lstlisting}

% ==============================================================================
\section{Lit and FileCheck Regression Testing Harness}
\label{sec:app_testing_harness}
% ==============================================================================

To verify optimization passes in continuous integration pipelines:

\begin{lstlisting}[language=Python, caption={LLVM Lit Test Suite Configuration (test/lit.cfg.py)}]
import os
import lit.formats

config.name = "DQC"
config.test_format = lit.formats.ShTest(True)
config.suffixes = [".mlir", ".ll"]

config.test_source_root = os.path.dirname(__file__)
config.test_exec_root = os.path.join(config.my_obj_root, "test")

config.substitutions.append(("%PATH%", config.environment["PATH"]))
\end{lstlisting}

\begin{lstlisting}[language=MLIR, caption={Sample FileCheck Test for Pass 3 (test/DQC/teleportation.mlir)}]
// RUN: dqc-opt %s -dqc-teleportation-synthesis | FileCheck %s

// CHECK-LABEL: func.func @test_remote_cnot
func.func @test_remote_cnot() {
  %q0 = dqc.alloc_qubit on 0 : !dqc.qubit
  %q1 = dqc.alloc_qubit on 1 : !dqc.qubit

  // CHECK: %epr = dqc.prepare_epr 0 <-> 1 fid 0.960
  // CHECK: %{{.*}}, %{{.*}} = dqc.telegate %{{.*}}, %{{.*}}, %epr {kind = "CNOT"}
  %0, %1 = dqc.cnot %q0, %q1 : !dqc.qubit, !dqc.qubit
  return
}
\end{lstlisting}

\begin{lstlisting}[language=MLIR, caption={Sample FileCheck Test for Pass 4 JIT Scheduling (test/DQC/scheduling.mlir)}]
// RUN: dqc-opt %s -dqc-coherence-scheduling | FileCheck %s

// CHECK-LABEL: func.func @test_jit_hoisting
func.func @test_jit_hoisting() {
  %q0 = dqc.alloc_qubit on 0 : !dqc.qubit
  %q1 = dqc.alloc_qubit on 1 : !dqc.qubit

  // CHECK: dqc.prepare_epr {{.*}} {sched.time_ns = 2000.0}
  %epr = dqc.prepare_epr 0 <-> 1 fid 0.98 : !dqc.epr_pair
  %0, %1 = dqc.telegate %q0, %q1, %epr {kind = "CNOT"} : !dqc.qubit, !dqc.qubit
  return
}
\end{lstlisting}

% ==============================================================================
\section{Python Bindings and MLIR C-API Integration}
\label{sec:app_python_bindings}
% ==============================================================================

To allow data scientists and quantum algorithms researchers to drive compilation from Python and Jupyter notebooks:

\begin{lstlisting}[language=C++, caption={pybind11 MLIR Python Bindings (python/dqc/DQCModule.cpp)}]
#include <pybind11/pybind11.h>
#include "mlir-c/Bindings/Python/Interop.h"
#include "dqc/DQCDialect.h"

namespace py = pybind11;

PYBIND11_MODULE(_dqcOps, m) {
  m.doc() = "DQC MLIR Dialect Python Bindings";

  m.def("register_dialect", [](MlirContext context) {
    MlirDialectHandle handle = mlirGetDialectHandle__dqc__();
    mlirDialectHandleRegisterDialect(handle, context);
  });
}
\end{lstlisting}

\begin{lstlisting}[language=Python, caption={Python SDK Script Executing DQC Passes (example.py)}]
import dqc
from mlir.ir import Context, Module

with Context() as ctx:
    dqc.register_dialect(ctx)
    mod = Module.parse("""
    module {
      func.func @bell() {
        %q0 = dqc.alloc_qubit on 0 : !dqc.qubit
        %q1 = dqc.alloc_qubit on 1 : !dqc.qubit
        %0 = dqc.h %q0 : !dqc.qubit
        %1, %2 = dqc.cnot %0, %q1 : !dqc.qubit, !dqc.qubit
        return
      }
    }
    """)
    
    # Run 5-Pass Pipeline
    pm = dqc.PassManager.from_pipeline("dqc-pipeline")
    pm.run(mod.operation)
    print(mod)
\end{lstlisting}
